% =============================================================================
% SVVVL Research Paper — Higher-Order Root Finding via Scale Ratios
% House-style preamble (shared); edit content, not the preamble.
% =============================================================================

\documentclass[11pt]{article}

% ----- Page geometry -------------------------------------------------------
\usepackage[
  letterpaper,
  top=0.9in,
  bottom=0.9in,
  left=0.85in,
  right=0.85in,
  headheight=14pt
]{geometry}

% ----- Encoding & language -------------------------------------------------
\usepackage[utf8]{inputenc}
\usepackage[T1]{fontenc}
\usepackage[english]{babel}



% ----- Typography ----------------------------------------------------------
\usepackage{XCharter}

% Section headings
\usepackage{titlesec}
\titleformat{\section}
  {\normalfont\Large\bfseries}{\thesection}{0.6em}{}
\titleformat{\subsection}
  {\normalfont\large\bfseries}{\thesubsection}{0.6em}{}
\titlespacing*{\section}{0pt}{1.4em}{0.6em}
\titlespacing*{\subsection}{0pt}{1.1em}{0.4em}

% Remove section numbering
\setcounter{secnumdepth}{0}

% ----- Spacing & paragraphs ------------------------------------------------
\usepackage{setspace}
\setstretch{1.15}
\setlength{\parskip}{0.6em}
\setlength{\parindent}{0pt}

% ----- Graphics & color ----------------------------------------------------
\usepackage{graphicx}
\usepackage{xcolor}
\definecolor{linkcolor}{HTML}{1F4E79}

% ----- Hyperlinks ----------------------------------------------------------
\usepackage[
  colorlinks=true,
  linkcolor=linkcolor,
  urlcolor=linkcolor,
  citecolor=linkcolor
]{hyperref}

% ----- Tables, figures, captions ------------------------------------------
\usepackage{booktabs}
\usepackage{array}
\usepackage{caption}
\captionsetup{font=small, labelfont=bf, skip=6pt}

% ----- Math ----------------------------------------------------------------
\usepackage{amsmath, amssymb}
\usepackage{amsthm}
\usepackage{mathtools}

% ----- Footnotes -----------------------------------------------------------
\usepackage[hang, flushmargin]{footmisc}
\setlength{\footnotesep}{0.8em}

% ----- Headers & footers ---------------------------------------------------
\usepackage{fancyhdr}
\pagestyle{fancy}
\fancyhf{}
\renewcommand{\headrulewidth}{0pt}
\fancyfoot[C]{\thepage}

\fancypagestyle{firstpage}{%
  \fancyhf{}%
  \renewcommand{\headrulewidth}{0pt}%
}


% ----- Logo lockup ---------------------------------------------------------
\newcommand{\WORDMARKFILE}{../../assets/SVVVL_Research_Wordmark.pdf}
\newcommand{\WORDMARKWIDTH}{1.5in}
\newcommand{\LOCKUPGAP}{0.01in}

% The lockup itself: symbol above wordmark, both centered within a box
% the width of the symbol so the wordmark sits naturally beneath it.
\newcommand{\logolockup}{%
  \begin{minipage}{\WORDMARKWIDTH}
    \centering
    \includegraphics[width=\WORDMARKWIDTH]{\WORDMARKFILE}
  \end{minipage}%
}

% Title block command: \makecover{Title}{Subtitle}{Authors}{Date}
\newcommand{\makecover}[4]{%
  \thispagestyle{firstpage}%
  \noindent\logolockup\par
  \vspace{2.2em}
  {\fontsize{26pt}{30pt}\selectfont\bfseries
   \hyphenpenalty=10000\exhyphenpenalty=10000\sloppy
   #1\par}
  \vspace{0.6em}
  {\fontsize{15pt}{18pt}\selectfont\itshape\color{black!70}
   \hyphenpenalty=10000\exhyphenpenalty=10000\sloppy
   #2\par}
  \vspace{1.6em}
  {\normalsize\color{black!60} #3\par}
  \vspace{0.2em}
  {\normalsize\color{black!60} #4\par}
  \vspace{1.8em}
}

% ----- Theorem environments ------------------------------------------------
\newtheorem{proposition}{Proposition}
\newtheorem{remark}{Remark}
\newtheorem{conjecture}{Conjecture}
\newtheorem{theorem}{Theorem}
\newtheorem{lemma}{Lemma}
\newtheorem{corollary}{Corollary}
\newtheorem{definition}{Definition}

% ----- Custom commands -----------------------------------------------------
% General-purpose math operators reused across papers.
\newcommand{\E}{\mathbb{E}}
\newcommand{\Var}{\mathrm{Var}}
\newcommand{\Corr}{\mathrm{Corr}}
\newcommand{\dd}{\mathrm{d}}
% Paper-specific notation: characteristic length scales and their ratios.
\newcommand{\Lf}{L_{f}}                 % leading length scale f/f'
\newcommand{\Lk}[1]{L_{#1}}             % L_k = f^{(k)}/f^{(k+1)}
\newcommand{\rr}{\rho}                  % scale ratio
\newcommand{\order}{d}                  % method order parameter

% =============================================================================
\begin{document}
% =============================================================================

% --- COVER -------------------------------------------------------------------
\makecover
  {Higher-Order Root Finding as a Hierarchy of Scales}
  {Newton, Halley, and the Householder family read off the characteristic
   length scales of a function}
  {J.~E. Whitehead~III\footnote{SVVVL, Towson, MD, United States.
   Contact: \href{mailto:james.whitehead@svvvl.com}{james.whitehead@svvvl.com}.
   See Disclosures at the end of this document.}}
  {\today}

% --- ABSTRACT ----------------------------------------------------------------
\begin{quote}
\itshape
Newton's method steps by $-f/f'$. Read $f/f'$ as a length, not as a ratio of
derivatives: it is the characteristic scale $\Lf$ of the exponential that matches
$f$ at the current point, so Newton steps one scale-length toward the root. Every
higher-order method is that same step, corrected for how the scale drifts as you
move. We make the correction explicit and read off three pictures of it. Define a
tower of scales $\Lk{k}=f^{(k)}/f^{(k+1)}$ and the ratios $\rr_k=\Lk{k-1}/\Lk{k}$
between rungs. Halley is Newton divided by $1-\tfrac12\rr_1$; the order-$\order$
Householder method is $\Lf$ times a fixed rational function of
$\rr_1,\dots,\rr_{\order-1}$ --- one continued fraction, read one rung deeper at
each order. The step is a clock, the time-to-contact of a particle closing on a
target; and because a root of $f$ is a pole of $1/f$, it is the ratio test ranging
to that pole. We close with a root finder that needs only the recursion for the
coefficients of $1/f$ and one ratio per step.
\end{quote}


\newpage
% --- INTRODUCTION ------------------------------------------------------------
\section{Newton's method is a statement about length}

Newton's method updates a guess $x_t$ for a root of $f$ by
\begin{equation}
  x_{t+1} = x_t - \frac{f(x_t)}{f'(x_t)}
          = x_t - \frac{1}{\,f'(x_t)/f(x_t)\,}.
  \label{eq:newton}
\end{equation}
The textbook reading is algebraic: solve the local linear model
$f(x)+f'(x)(y-x)=0$ for its zero. Here is a different reading, the one that
generalizes. The ratio $f'/f$ is the logarithmic derivative, and it has units of
inverse length. So its reciprocal,
\begin{equation}
  \Lf(x) \;\coloneqq\; \frac{f(x)}{f'(x)},
  \label{eq:Lf}
\end{equation}
is a length --- and it is \emph{the} length built into $f$ at $x$. Fit a local
exponential $f(y)\approx A\,e^{y/L}$ to the value and slope of $f$ at $x$, and the
fit hands back $L=\Lf(x)$, because $g/g'=L$ everywhere for $g=A\,e^{y/L}$. So $\Lf$
is the distance over which $f$ changes by a factor of $e$, and Newton's
rule~\eqref{eq:newton} reads:

\begin{quote}
\itshape
Step one characteristic length toward the root: $x_{t+1}=x_t-\Lf(x_t)$.
\end{quote}

That is more than a mnemonic. The exponential is the one function whose scale never
changes, and it is exactly the function on which a single scale-length is the right
unit of progress. Real functions have a scale that drifts as $x$ moves, and that
drift is what Newton's first-order step gets wrong. Every higher-order method is
bookkeeping for the drift.

Here is the plan. Section~2 makes the drift quantitative. Section~3 shows Halley is
Newton with the first drift correction. Section~4 reads the step as a clock --- the
physical heart of the paper. Section~5 writes the whole Householder family as one
continued fraction in the scale ratios. Section~6 turns the integer order into a
continuous dial and finds the root sitting where $1/f$ blows up. Section~7 folds all
of it into an algorithm you can run. Don't worry if the length picture feels loose
now; each section sharpens it.

% --- THE TOWER OF SCALES -----------------------------------------------------
\section{The scale drifts, and that drift is what Newton misses}

If $\Lf=f/f'$ is the scale of $f$, then $f'$ has its own scale $f'/f''$, and $f''$
has $f''/f'''$, on up. Define the tower
\begin{equation}
  \Lk{k}(x) \;\coloneqq\; \frac{f^{(k)}(x)}{f^{(k+1)}(x)},
  \qquad k=0,1,2,\dots,
  \label{eq:tower}
\end{equation}
so $\Lk{0}=\Lf$ is the leading scale and each $\Lk{k}$ is the characteristic length
of the $k$-th derivative. The same information as the derivatives $f,f',f'',\dots$,
carried in the currency of lengths.

A single length can't tell you how the scale is changing; for that you need the
dimensionless ratio between neighboring rungs,
\begin{equation}
  \rr_k(x) \;\coloneqq\; \frac{\Lk{k-1}(x)}{\Lk{k}(x)},
  \qquad k=1,2,\dots.
  \label{eq:ratio}
\end{equation}
The first one is worth seeing in full:
\begin{equation}
  \rr_1 \;=\; \frac{\Lk{0}}{\Lk{1}}
        \;=\; \frac{f/f'}{f'/f''}
        \;=\; \frac{f\,f''}{(f')^{2}}.
  \label{eq:rho1}
\end{equation}
Two reference functions fix its meaning. For a \emph{line} $f=a(x-x^\star)$ we have
$f''=0$, so $\rr_1=0$: the scale isn't bending. For an \emph{exponential}
$f=e^{x/L}$ every rung equals $L$ and every $\rr_k=1$: the scale is perfectly
self-similar. In between, $\rr_1$ measures how far $f$ has bent from a straight line
toward an exponential.

Now the one quantity Newton ignores. Differentiate the leading scale~\eqref{eq:Lf}:
\begin{equation}
  \Lf'(x)
   = \frac{(f')^2 - f f''}{(f')^2}
   = 1 - \frac{f f''}{(f')^2}
   = 1 - \rr_1.
  \label{eq:drift}
\end{equation}
\textbf{The leading scale drifts at rate $\Lf'=1-\rr_1$.} In plain terms: Newton
steps by $-\Lf$ as if $\Lf$ were the fixed distance to the root, but $\Lf$ moves as
$x$ moves, and $\rr_1$ says how fast. A line has $\Lf'=1$ --- the scale equals the
distance to the root and runs out right at it. An exponential has $\Lf'=0$ --- the
scale never moves. So the first correction any second-order method makes can depend
on nothing but $\rr_1$. That points straight at Halley.

% --- HALLEY -----------------------------------------------------------------
\section{Halley repairs the drift}

Halley's second-order method, in these variables, is Newton with one correction
factor:
\begin{equation}
  x_{t+1} = x_t - \frac{\Lf}{\,1 - \tfrac12\,\dfrac{\Lf}{\Lk{1}}\,}
          = x_t - \frac{\Lf}{\,1 - \tfrac12\,\rr_1\,}.
  \label{eq:halley}
\end{equation}
This is the standard Halley iteration~[\hyperlink{ref:4}{4},\,\hyperlink{ref:5}{5}].
To see it, take the textbook form and divide top and bottom by $2(f')^2$:
\begin{equation}
  x_{t+1}
   = x_t - \frac{2 f f'}{2 (f')^2 - f f''}
   = x_t - \frac{f/f'}{\,1 - \dfrac{f f''}{2 (f')^2}\,}
   = x_t - \frac{\Lf}{\,1 - \tfrac12\rr_1\,},
  \label{eq:halley-derive}
\end{equation}
using~\eqref{eq:rho1} for the last step. Halley keeps Newton's direction --- the
leading scale $\Lf$ --- and rescales the step by $1/(1-\tfrac12\rr_1)$, a factor
built entirely from the first scale ratio. When $\rr_1>0$ (the scale shrinking ahead
of the step) it lengthens Newton's step; when $\rr_1<0$ it shortens it.

\textbf{Halley is Newton plus a first-order correction for the drift.} It accounts
for how the length scale changes under the very step it is about to take. Before
climbing to higher orders, look at what this step \emph{is}, physically.

% --- KINEMATICS -------------------------------------------------------------
\section{The kinematic reading: a clock counting down to contact}

Why does $-f/f'$ keep showing up? Because it is a time.

Picture $f$ along $x$ as the signed miss-distance of a particle closing on a target,
with each step of $x$ playing the role of elapsed time. Then the local derivatives
are the particle's state of motion: $f$ is position, $f'$ is velocity --- the closing
rate --- $f''$ is acceleration, $f'''$ is jerk, and on up the tower.

In those terms the characteristic length is a clock. Distance over closing rate is
the first-order time to contact,
\begin{equation}
  \tau = -\frac{f}{f'} = -\Lf,
  \label{eq:ttc}
\end{equation}
the same $-\Lf$ Newton steps by --- and the same quantity that visual guidance calls
``$\tau$'' and missile guidance calls time-to-go~[\hyperlink{ref:7}{7}].

\begin{quote}\itshape
Newton steps by the time-to-contact at constant closing speed.
\end{quote}

The drift law says this clock keeps honest time. Since $\Lf'=1-\rr_1$, the
time-to-contact changes with position at rate $\dd\tau/\dd x=\rr_1-1$, which is
exactly $-1$ when the motion is uniform ($\rr_1=0$): advance one step, lose one tick.
The dimensionless acceleration $\rr_1=f f''/(f')^2$ is the correction when the closing
rate is not constant, and Halley~\eqref{eq:halley} folds it in:
$\tau=\tau_1/(1-\tfrac12\rr_1)$ is the time-to-contact corrected for acceleration.
Each higher order adds the next term of the motion --- order 3 the jerk, order 4 the
snap --- so the order-$\order$ method reads the contact time off the kinematics
through the $\order$-th derivative. The next section gives that family a single form.

% --- THE HOUSEHOLDER FAMILY -------------------------------------------------
\section{The whole family is one continued fraction}

Newton and Halley are the first two rungs of one family. The order-$\order$
Householder method is
\begin{equation}
  x_{t+1} = x_t + \order\,
            \frac{\big(1/f\big)^{(\order-1)}}{\big(1/f\big)^{(\order)}},
  \label{eq:householder}
\end{equation}
with $\order=1$ Newton and $\order=2$ Halley. Each member is a ratio of consecutive
derivatives of $1/f$, and dividing through by the leading scale collapses all of that
into the dimensionless ratios: the order-$\order$ step is $-\Lf$ times a fixed
rational function of $\rr_1,\dots,\rr_{\order-1}$. The first three:
\begin{align}
  \text{Newton } (\order=1):\quad
    & x_{t+1} = x_t - \Lf, \\[0.3em]
  \text{Halley } (\order=2):\quad
    & x_{t+1} = x_t - \Lf\,\frac{1}{\,1-\tfrac12\rr_1\,}, \\[0.3em]
  \text{order } 3:\quad
    & x_{t+1} = x_t - \Lf\,
      \frac{1-\tfrac12\rr_1}{\,1-\rr_1+\tfrac16\rr_1^{2}\rr_2\,}.
  \label{eq:order3}
\end{align}
The order-3 denominator is the one new piece; the numerator just repeats Halley.
Carry the computation higher --- by hand it turns ugly, so we ran it in a
computer-algebra system through order 5 --- and a clean pattern shows up. Define a
sequence of polynomials in the scale ratios,
\begin{align}
  P_0 = P_1 &= 1, \notag\\
  P_2 &= 1 - \tfrac12\rr_1, \notag\\
  P_3 &= 1 - \rr_1 + \tfrac16\rr_1^2\rr_2.
  \label{eq:Ppolys}
\end{align}
In terms of these, the order-$\order$ step is just
\begin{equation}
  x_{t+1} = x_t - \Lf\,\frac{P_{\order-1}}{P_\order}.
  \label{eq:hhstep}
\end{equation}
\textbf{The numerator at each order is the denominator of the order before.} So the
family is one continued fraction in the scale ratios, read one rung deeper at each
order --- not a list of separate tricks. Each order folds in exactly one new ratio,
$\rr_{\order-1}$, riding the first scale that needs the $\order$-th derivative.\footnote{%
The whole table closes to one formula, $P_\order=(-\Lf)^{\order}\,f\,(1/f)^{(\order)}/\order!$,
which is just ``take one more derivative of $1/f$'' said in the scale variables. Its
coefficients are a sum over the partitions of $\order$ by Fa\`a di Bruno's
formula~[\hyperlink{ref:6}{6}]; we verified the telescoping through order~7.}

One case makes the structure vivid. On a pure exponential every rung equals $L$,
every $\rr_k=1$, and $P_\order=1/\order!$, so~\eqref{eq:hhstep} gives a step of
$-\Lf\,(1/(\order-1)!)\big/(1/\order!)=-\order\,\Lf$. Each order buys one more
e-folding of progress on the scale-invariant function. That raises the obvious
question: turn the order up without limit, and where does the step land?

% --- POLE / CONTINUUM -------------------------------------------------------
\section{A root is a pole, and the order is a dial you can turn}

To answer it, look at what the step actually computes. Write the Taylor coefficients
of $f$ and of $1/f$ about $x$ as $f_n=f^{(n)}(x)/n!$ and $g_n=(1/f)^{(n)}(x)/n!$. Then
the order-$\order$ step~\eqref{eq:householder} is exactly the ratio of two consecutive
coefficients of $1/f$,
\begin{equation}
  \mathrm{step}_\order
   = \order\,\frac{(1/f)^{(\order-1)}}{(1/f)^{(\order)}}
   = \frac{g_{\order-1}}{g_\order}.
  \label{eq:ratiostep}
\end{equation}
And those coefficients come free from one identity, $f\cdot(1/f)=1$, read off order by
order:
\begin{equation}
  \sum_{k=0}^{n} f_k\,g_{n-k} = \delta_{n,0}.
  \label{eq:convorth}
\end{equation}
This solves from the top down for each $g_n$ in terms of the coefficients already in
hand --- the recursion the algorithm in Section~7 will use.

Now the picture. A root of $f$ is a pole of $1/f$.

\begin{quote}\itshape
The Householder step is the ratio test ranging to that pole.
\end{quote}

Here is why, made concrete. Near a simple root $x^\star$, the reciprocal $1/f$ has a
simple pole, and its coefficients about $x$ grow like
$g_n\propto(x^\star-x)^{-(n+1)}$. Their ratio is the distance itself,
$g_{\order-1}/g_\order\to x^\star-x$ --- the ratio test for the radius of
convergence, which is the distance to the nearest singularity. With more than one
root, the step is a moment-ratio estimate of the nearest, sharpening as the order
climbs onto higher coefficients. Table~\ref{tab:ratiotest} shows it homing in for
$f=(x-3)(x-7)$ started at $x=0$.

\begin{table}[ht]
\centering
\caption{The step $g_{\order-1}/g_\order$ for $f=(x-3)(x-7)$ started at $x=0$, homing
onto the nearest root $x^\star=3$. Newton ($\order=1$) undershoots because the far
root at $7$ still pulls; each higher order leans on higher coefficients of $1/f$ and
ranges the near root more sharply.}
\label{tab:ratiotest}
\renewcommand{\arraystretch}{1.3}
\begin{tabular}{l c c c c c c c}
\toprule
order $\order$ & 1 & 2 & 3 & 4 & 5 & 6 & 7 \\
\midrule
$g_{\order-1}/g_\order$ & 2.100 & 2.658 & 2.860 & 2.941 & 2.975 & 2.989 & 2.995 \\
\bottomrule
\end{tabular}
\end{table}

That convergence has a closed-form limit, and finding it is where the continuous
order comes in. The only discrete thing in the whole construction is the factorial
$\order!$. And the factorial is already an integral,
\begin{equation}
  \order! = \Gamma(\order+1) = \int_0^\infty u^{\order}\,e^{-u}\,\dd u,
  \label{eq:gamma}
\end{equation}
the $\order$-th moment of a unit-rate exponential. Replace $\order!$ by
$\Gamma(\order+1)$ and the order becomes a knob you can turn to any real value, even
to infinity.

Turn it. Near the root, write $\lambda=x^\star-x$ for the distance. The coefficients
$g_n=\lambda^{-(n+1)}$ are the normalized moments of an exponential density of rate
$\lambda$, and summing them rebuilds $1/f$ as that density's Laplace
transform~[\hyperlink{ref:8}{8}],
\begin{equation}
  \sum_{n\ge0} g_n\,s^{n}
   = \int_0^\infty e^{-\lambda u}\,e^{s u}\,\dd u
   = \frac{1}{\lambda-s}
   = \frac{1}{f(x+s)}.
  \label{eq:mgf}
\end{equation}
Push the order to any real $\nu$ and the gamma cancels, leaving the step equal to the
distance at every order:
\begin{equation}
  g_\nu = \frac{1}{\Gamma(\nu+1)}\int_0^\infty u^{\nu}\,e^{-\lambda u}\,\dd u
        = \lambda^{-(\nu+1)},
  \qquad
  \frac{g_{\nu-1}}{g_\nu} = \lambda = x^\star-x.
  \label{eq:contstep}
\end{equation}
\textbf{On a pure pole the method is exact the moment you write it down.}

A general $f$ is not a single pole. But for analytic $f$ the reciprocal $1/f$ is
meromorphic --- a pole at every root --- and its partial-fraction expansion writes it
as the Laplace transform of one exponential per root, each decaying at the rate of its
own distance $\lambda_j=x^\star_j-x$ and weighted by the residue $w_j$ of $1/f$ there:
\begin{equation}
  \frac{1}{f(x+s)} = \sum_j \frac{w_j}{\lambda_j-s}
   = \int_0^\infty e^{su}\sum_j w_j\,e^{-\lambda_j u}\,\dd u,
  \qquad
  g_n = \sum_j w_j\,\lambda_j^{-(n+1)}.
  \label{eq:spectral}
\end{equation}
Each coefficient is a residue-weighted blend of modes, and the step leans on the
slowest --- the nearest root --- ever harder as the order climbs.\footnote{This needs
analyticity, not mere smoothness: a $C^\infty$ but non-analytic $f$ has no convergent
local series. The spectral weights are in general signed, even complex, not a
probability density, and the one-sided transform is cleanest for roots ahead of $x$.
For $f=(x-3)(x-7)$ about $x=0$ the two modes are $\tfrac14 e^{-3u}$ and
$-\tfrac14 e^{-7u}$, and the near one wins (Table~\ref{tab:ratiotest}). A root of
multiplicity $m$ contributes an Erlang mode $u^{m-1}e^{-\lambda u}$; a
complex-conjugate pair, a damped oscillation.} That is why the integer steps in the
Table climb toward $\lambda$: $2.100\to2.995$ onto the root at $x^\star=3$, now with a
limit you can name.

This closes the circle. The exponential was the one function with a single fixed scale
and no drift --- every $\rr_k=1$, the fixed point of Section~2. It returns here as the
\emph{law} whose Laplace transform is $1/f$ near the root: the same function at both
ends, once as the shape Newton steps along, once as the distribution the step reads.

% --- ALGORITHM --------------------------------------------------------------
\section{An algorithm: step by the ratio test}

Everything above turns into a short loop. The order-$\order$ Householder method never
needs the symbolic $\order$-th derivative of $1/f$; it needs only the coefficients
$g_n$, and those come from the recursion~\eqref{eq:convorth} in $O(\order^2)$ work. The
step is one ratio.

\begin{quote}
\ttfamily
\textbf{input} $x$, order $\order$, tolerance $\varepsilon$, cap $N$\\[0.3em]
\textbf{repeat}\\
\hspace*{1.5em}$f_k \gets f^{(k)}(x)/k!$ \quad for $k=0,\dots,\order$\\
\hspace*{1.5em}$g_0 \gets 1/f_0$\\
\hspace*{1.5em}$g_n \gets -\,f_0^{-1}\textstyle\sum_{k=1}^{n} f_k\,g_{n-k}$ \quad for $n=1,\dots,\order$\\
\hspace*{1.5em}$\Delta \gets g_{\order-1}\,/\,g_\order$\\
\hspace*{1.5em}$x \gets x + \Delta$\\
\textbf{until} $|\Delta| < \varepsilon$ \ \textbf{or} \ iterations $> N$\\[0.3em]
\textbf{return} $x$
\end{quote}

Three notes make it practical. \emph{Firstly}, the cost per iteration is the $\order$
derivatives of $f$ plus an $O(\order^2)$ convolution --- cheap, and the same
derivatives any high-order method would want. \emph{Secondly}, the order $\order$ is
yours to set, and it buys convergence speed:

\begin{proposition}[Convergence order]
If $x^\star$ is a simple root of a sufficiently smooth $f$, the order-$\order$
iteration~\eqref{eq:householder} converges locally with order $\order+1$: Newton
quadratically, Halley cubically, order 3 quartically.
\label{prop:conv}
\end{proposition}

We don't reprove this classical fact; see Householder~[\hyperlink{ref:1}{1}] and
Traub~[\hyperlink{ref:2}{2}]. The structural reason is the one the paper has been
making: the order-$\order$ step uses the scales $\Lk{0},\dots,\Lk{\order-1}$ ---
derivatives through $f^{(\order)}$ --- and cancels the error in the
distance-to-root expansion through that order. One more scale, one more ratio, one
more order of convergence; and, on the exponential, one more e-folding.

\emph{Thirdly}, the coefficients carry their own diagnostic. The ratios
$g_0/g_1,\,g_1/g_2,\,\dots,\,g_{\order-1}/g_\order$ --- the row of
Table~\ref{tab:ratiotest} --- should settle onto a single value as $n$ grows, and that
value is the distance to the nearest root. How fast they settle tells you whether the
order is high enough to trust the step; when they have not settled, the nearest root
has competition, and raising $\order$ sharpens the estimate.

\textbf{Build the coefficients of $1/f$, take their last ratio, step. That is the
whole method.}

% --- CONCLUSION --------------------------------------------------------------
\section{Conclusion}

Read $f/f'$ as a length and the classical root finders line up as one idea. Newton
steps one characteristic scale toward the root. The scale drifts at rate $1-\rr_1$,
and Halley divides by $1-\tfrac12\rr_1$ to correct for it. The order-$\order$
Householder method is the same leading scale $\Lf$ rescaled by a fixed rational
function of the consecutive ratios $\rr_1,\dots,\rr_{\order-1}$ --- one continued
fraction, one rung deeper per order. That single step wears three faces: a clock
reading the time-to-contact of a particle closing on a target; on a pure exponential
the clean $-\order\,\Lf$, one e-folding per order; and, because a root of $f$ is a
pole of $1/f$, the ratio test ranging to that pole. Turn that last reading into a loop
--- build the coefficients of $1/f$, take their last ratio, step --- and you have a
root finder whose order is a dial. Higher-order root finding is reading the hierarchy
of a function's own length scales one rung deeper.

\newpage
% --- REFERENCES --------------------------------------------------------------
\section{Further Reading and References}

\begin{itemize}\setlength{\itemsep}{0.4em}

  \item[{[1]}]\hypertarget{ref:1}{} Householder, A.~S. (1970).
    \textit{The Numerical Treatment of a Single Nonlinear Equation.}
    McGraw-Hill, New York.

  \item[{[2]}]\hypertarget{ref:2}{} Traub, J.~F. (1964).
    \textit{Iterative Methods for the Solution of Equations.}
    Prentice-Hall, Englewood Cliffs, NJ.

  \item[{[3]}]\hypertarget{ref:3}{} Ostrowski, A.~M. (1966).
    \textit{Solution of Equations and Systems of Equations.}
    2nd ed., Academic Press, New York.

  \item[{[4]}]\hypertarget{ref:4}{} Gander, W. (1985).
    ``On Halley's Iteration Method.''
    \textit{American Mathematical Monthly}, 92(2), 131--134.

  \item[{[5]}]\hypertarget{ref:5}{} Halley, E. (1694).
    ``Methodus nova accurata \& facilis inveniendi radices
    \ae quationum quarumcumque generaliter, sine praevia reductione.''
    \textit{Philosophical Transactions of the Royal Society of London},
    18, 136--148.

  \item[{[6]}]\hypertarget{ref:6}{} Comtet, L. (1974).
    \textit{Advanced Combinatorics: The Art of Finite and Infinite Expansions.}
    Reidel, Dordrecht. (Fa\`a di Bruno's formula and Bell polynomials.)

  \item[{[7]}]\hypertarget{ref:7}{} Lee, D.~N. (1976).
    ``A theory of visual control of braking based on information about
    time-to-collision.'' \textit{Perception}, 5(4), 437--459.

  \item[{[8]}]\hypertarget{ref:8}{} Widder, D.~V. (1941).
    \textit{The Laplace Transform.} Princeton University Press, Princeton, NJ.

\end{itemize}

\newpage
% --- DISCLOSURES -------------------------------------------------------------
\section{Disclosures}

The views expressed in this document are those of the author and
SVVVL alone and do not constitute investment, legal, tax, or
accounting advice. This document is provided for informational and
research purposes only and is not an offer or solicitation to buy or
sell any security or financial product.

The framework presented is theoretical and relies on idealized
assumptions that do not hold in practice. The results
should not be implemented in a real portfolio without additional
constraints and considerations not addressed in this paper.

While reasonable care has been taken in the preparation of this
document, no representation or warranty, express or implied, is made
as to its accuracy or completeness.

Reading this document does not create an advisory, fiduciary, or
client relationship of any kind between the reader and SVVVL or the
author.

\end{document}
